% Unidad U014. Biografía estructural completa del carácter de autoescala.

\label{gfp:b:chap:biografia-completa-phi}

\section{Selección finita de la semilla de autoescala}

Sea \(\mathcal U_6\subset\{0,\ldots,999\}^6\) el catálogo de las
\(468\) emisiones distintas construido en el
\cref{gfp:b:chap:semillas-emisor-54}, y sea
\[
 \Psi:\mathcal U_6\longrightarrow\mathbb F_3^6,
 \qquad
 \Psi(U)=((-U_j)\bmod3)_{j=1}^{6}.
\]
Para \(w\in\Psi(\mathcal U_6)\), recordamos los invariantes
\[
 n(w)=\#\Psi^{-1}(w),
 \qquad
 M(w)=\sum_{U\in\Psi^{-1}(w)}\operatorname{mult}(U),
\]
y el vector de ocupación
\[
 \nu(w)=
 \bigl(\#\{j:w_j=0\},\#\{j:w_j=1\},\#\{j:w_j=2\}\bigr).
\]

\begin{theorem}[Selección orbital de la autoescala]
\label{gfp:b:thm:seleccion-orbital-phi}
La acción diedral sobre \(\Psi(\mathcal U_6)\) contiene una única órbita
que satisface el predicado de autoescala
\[
 n(w)=1,
 \qquad
 M(w)=144,
 \qquad
 \nu(w)=(2,2,2).
\]
El calibre interno
\[
 \operatorname{diag}_c=6,
 \qquad
 \operatorname{card}=ES
\]
selecciona en esa órbita la palabra
\[
 \boxed{w_6^{\varphi}=121200.}
\]
Su única emisión es
\[
 \boxed{U_{\varphi}=(830,943,830,109,252,252),}
\]
de multiplicidad \(144\), firma multiplicativa \(555933\), orientación
aditiva \(ES\), soporte multiplicativo de cardinal seis y tipo residual
\(A\).
\end{theorem}

\begin{proof}
La proyección directa da
\[
 U_{\varphi}\bmod3=(2,1,2,1,0,0),
 \qquad
 -U_{\varphi}\bmod3=(1,2,1,2,0,0).
\]
La enumeración se efectúa sobre las \(243\) palabras visibles y sus
\(43\) órbitas, no sobre una selección de ejemplos. El predicado conjunto
\(n=1\), \(M=144\), \(\nu=(2,2,2)\) deja una órbita. Dentro de ella, la
diagonal y la orientación declaradas dejan una emisión. En esta etapa no
se ha introducido la ecuación \(x^2-x-1=0\), una tabla decimal ni el nombre
convencional de la constante.
\end{proof}

\section{Elevación por bloques y \(1^{\mathrm a}\) frontera coinductiva}

Sea \(b_1=w_6^{\varphi}\), y sea
\(\varepsilon=(0,0,1,1)\). Con las matrices \(L_0,L_1\) construidas en el
\cref{gfp:b:chap:orbitas-elevacion-r36}, definimos
\[
 b_{k+1}=b_kL_{\varepsilon_k}\pmod3
 \qquad(1\leq k\leq4),
\]
y concatenamos
\[
 w_{6(k+1)}^{\varphi}=w_{6k}^{\varphi}\mathbin{\|}b_{k+1}.
\]
El cálculo matricial produce, bloque por bloque,
\[
\begin{aligned}
 b_1&=121200,\\
 b_2&=112202,\\
 b_3&=121020,\\
 b_4&=010210,\\
 b_5&=010200,
\end{aligned}
\]
y, por tanto,
\begin{equation}
 \boxed{
 w_{30}^{\varphi}
 =121200|112202|121020|010210|010200.}
 \label{gfp:b:eq:w30-phi}
\end{equation}
La neutralidad dual de la frontera añade
\[
 u_5^{\varphi}=102011,
\]
de modo que
\begin{equation}
 \boxed{
 w_{36}^{\varphi}
 =121200|112202|121020|010210|010200|102011.}
 \label{gfp:b:eq:w36-phi}
\end{equation}

\begin{remark}[Dos usos distintos de una palabra visible]
La palabra \(102011\) es también visible en la biografía de \(e\). La
igualdad de una proyección no identifica los estados enriquecidos: el canal,
el prefijo, el cociente, el acarreo, la fase y la memoria pertenecen al tipo
del objeto. Esta separación es obligatoria antes de construir cualquier
carácter real.
\end{remark}

\section{Del calendario trítico al multigrafo de incidencia}

\subsection{Regla modal primitiva}

El calendario trítico primitivo es
\[
 (+1,-1,0).
\]
Fijamos dos modos visibles, \(\Sigma\) y \(\Pi\), y la regla
\begin{equation}
 +1:m\longmapsto\Sigma,
 \qquad
 -1:m\longmapsto\Pi,
 \qquad
 0:m\longmapsto m.
 \label{gfp:b:eq:regla-modal-primitiva-phi}
\end{equation}
Con la condición de cierre cíclico, la órbita modal producida por
\eqref{gfp:b:eq:regla-modal-primitiva-phi} es
\[
 \boxed{(\Sigma,\Pi,\Pi).}
\]
Los tres pares consecutivos son
\[
 \Sigma\longrightarrow\Pi,
 \qquad
 \Pi\longrightarrow\Pi,
 \qquad
 \Pi\longrightarrow\Sigma.
\]

\begin{definition}[Matriz de incidencia cronológica]
Para \(i,j\in\{\Sigma,\Pi\}\), sea
\[
 N_3(i,j)=
 \#\{r\in\mathbb Z/3\mathbb Z:(m_r,m_{r+1})=(i,j)\}.
\]
El subíndice \(3\) cuenta instantes del calendario primitivo; no denota una
potencia matricial.
\end{definition}

\begin{theorem}[Descenso necesario de la incidencia modal]
\label{gfp:b:thm:descenso-incidencia-modal-phi}
En el orden \((\Sigma,\Pi)\),
\begin{equation}
 \boxed{
 N_3=
 \begin{pmatrix}
 0&1\\
 1&1
 \end{pmatrix}.}
 \label{gfp:b:eq:N3-phi}
\end{equation}
Por repetición exacta del calendario,
\begin{equation}
 \boxed{
 N_9=3N_3,
 \qquad
 N_{27}=9N_3,
 \qquad
 N_{108}=36N_3.}
 \label{gfp:b:eq:incidencias-9-27-108-phi}
\end{equation}
Estas tres igualdades son igualdades de conteos de aristas; no afirman
\(N_9=N_3^3\), \(N_{27}=N_3^9\) ni \(N_{108}=N_3^{36}\).
\end{theorem}

\begin{proof}
No aparece \(\Sigma\to\Sigma\); cada uno de los otros tres pares aparece
una vez. Esto determina las cuatro entradas de \(N_3\). Los calendarios de
longitud \(9\), \(27\) y \(108\) contienen respectivamente \(3\), \(9\)
y \(36\) copias del bloque primitivo. Cada repetición multiplica todos los
conteos por el mismo entero, lo que prueba
\eqref{gfp:b:eq:incidencias-9-27-108-phi}.
\end{proof}

\subsection{Publicación areal--volumétrica}

La etiqueta modal y la hoja geométrica son objetos distintos. La carta
de publicación es
\[
 \mathfrak g_{\rm av}(\Sigma)=\mathscr H_{\rm ar},
 \qquad
 \mathfrak g_{\rm av}(\Pi)=\mathscr H_{\rm vol}.
\]
Transportando \eqref{gfp:b:eq:N3-phi} por \(\mathfrak g_{\rm av}\), obtenemos
el operador de incidencia
\begin{equation}
 \boxed{
 F_{\rm av}=
 \begin{pmatrix}
 0&1\\
 1&1
 \end{pmatrix}.}
 \label{gfp:b:eq:Fav-derivada-phi}
\end{equation}
Así, las dos transiciones cruzadas, la única autorretención volumétrica y la
ausencia de autorretención areal se derivan del calendario; no son cuatro
entradas independientes.

\begin{remark}[Alcance del descenso]
Olvidar la fase elemental no produce por sí solo un cociente de Markov
fuerte del Topological Phase Kernel, porque el modo visible no determina
cuál de los tres operadores internos actuará en el siguiente paso. Lo que
sí desciende sin hipótesis adicional es el multigrafo de aristas de un
bloque completo. Su envolvente mínimo de memoria uno tiene matriz de
adyacencia \(F_{\rm av}\).
\end{remark}

\section{Rayo positivo invariante y unicidad de la autoescala}

\begin{theorem}[Rayo de Perron derivado]
\label{gfp:b:thm:rayo-perron-phi}
El cono positivo de \(F_{\rm av}\) contiene un único rayo propio. Si lo
normalizamos como \(v=(1,x)^{\mathsf T}\), entonces
\[
 x^2-x-1=0,
 \qquad
 x>0,
\]
y, por tanto,
\begin{equation}
 \boxed{x=\frac{1+\sqrt5}{2}=\varphi.}
 \label{gfp:b:eq:phi-perron}
\end{equation}
\end{theorem}

\begin{proof}
La ecuación \(F_{\rm av}v=\lambda v\) da, coordenada a coordenada,
\[
 x=\lambda,
 \qquad
 1+x=\lambda x.
\]
Al eliminar \(\lambda\) resulta \(x^2-x-1=0\). Sus raíces son
\((1\pm\sqrt5)/2\), y sólo la positiva pertenece al interior del cono. La
matriz \(F_{\rm av}\) es primitiva, pues \(F_{\rm av}^2\) tiene todas sus
entradas positivas; el teorema de Perron--Frobenius garantiza que ese rayo
positivo es simple y único. El nombre \(\varphi\) se asigna después de esta
construcción.
\end{proof}

\begin{corollary}[Ecuación escalar de autoescala]
\[
 \varphi=1+\frac1\varphi,
 \qquad
 \varphi^2=\varphi+1,
 \qquad
 \varphi^{-1}=\varphi-1.
\]
Estas identidades son consecuencias de \eqref{gfp:b:eq:phi-perron}; no seleccionan
la semilla ni la matriz.
\end{corollary}

\section{Publicación de la autoescala \(\varphi\) mediante potencias de Fibonacci y convergencia diofántica}

Sea \(F_0=0\), \(F_1=1\) y
\[
 F_{n+1}=F_n+F_{n-1}\qquad(n\geq1).
\]

\begin{theorem}[Potencias exactas de la incidencia]
\label{gfp:b:thm:potencias-fibonacci-phi}
Para todo \(n\geq1\),
\begin{equation}
 \boxed{
 F_{\rm av}^{\,n}=
 \begin{pmatrix}
 F_{n-1}&F_n\\
 F_n&F_{n+1}
 \end{pmatrix}.}
 \label{gfp:b:eq:potencias-Fav-fibonacci}
\end{equation}
\end{theorem}

\begin{proof}
Para \(n=1\), la identidad es la definición de \(F_{\rm av}\). Si vale
para \(n\), la multiplicación a la derecha por \(F_{\rm av}\) da
\[
 \begin{pmatrix}
 F_n&F_{n-1}+F_n\\
 F_{n+1}&F_n+F_{n+1}
 \end{pmatrix}
 =
 \begin{pmatrix}
 F_n&F_{n+1}\\
 F_{n+1}&F_{n+2}
 \end{pmatrix},
\]
que es el caso \(n+1\).
\end{proof}

\begin{theorem}[Identidad de Cassini]
\label{gfp:b:thm:cassini-phi}
Para todo \(n\geq1\),
\begin{equation}
 \boxed{F_{n+1}F_{n-1}-F_n^2=(-1)^n.}
 \label{gfp:b:eq:cassini-phi}
\end{equation}
\end{theorem}

\begin{proof}
El determinante de \(F_{\rm av}\) es \(-1\). Tomando determinantes en
\eqref{gfp:b:eq:potencias-Fav-fibonacci},
\[
 F_{n-1}F_{n+1}-F_n^2
 =\det(F_{\rm av}^{\,n})=(-1)^n.
\]
\end{proof}

Definimos \(r_n=F_{n+1}/F_n\) para \(n\geq1\). De Cassini se sigue
\[
 r_n-r_{n-1}
 =\frac{F_{n+1}F_{n-1}-F_n^2}{F_nF_{n-1}}
 =\frac{(-1)^n}{F_nF_{n-1}}.
\]
Por tanto, las aproximaciones alternan alrededor del rayo de Perron.

\begin{proposition}[Horquillas racionales de Perron]
\label{gfp:b:prop:horquillas-perron-phi}
Para \(m\geq1\),
\begin{equation}
 \frac{F_{2m+2}}{F_{2m+1}}
 <\varphi<
 \frac{F_{2m+1}}{F_{2m}},
 \label{gfp:b:eq:horquilla-fibonacci-phi}
\end{equation}
y la anchura es
\begin{equation}
 \boxed{
 \frac{F_{2m+1}}{F_{2m}}
 -\frac{F_{2m+2}}{F_{2m+1}}
 =\frac{1}{F_{2m}F_{2m+1}}.}
 \label{gfp:b:eq:anchura-fibonacci-phi}
\end{equation}
\end{proposition}

\begin{proof}
La alternancia anterior sitúa los cocientes pares por debajo y los impares
por encima de \(\varphi\). La fórmula de la anchura es
\eqref{gfp:b:eq:cassini-phi} con \(n=2m+1\), tras poner ambas fracciones sobre
denominador común. Como \(F_n\to\infty\), la anchura tiende a cero.
\end{proof}

\section{Incidencia de memoria de orden \(1\) y medida de Parry del envolvente simbólico}

Sea \(X_{\rm av}\) el menor subshift de tipo finito sobre
\(\{\mathscr H_{\rm ar},\mathscr H_{\rm vol}\}\) que admite exactamente
los tres pares observados en el calendario modal. Su matriz de adyacencia es
\(F_{\rm av}\).

\begin{theorem}[Transición y medida de Parry]
\label{gfp:b:thm:parry-phi}
Sea \(r=(1,\varphi)^{\mathsf T}\). La matriz estocástica de Parry es
\begin{equation}
 \boxed{
 P_{ij}=
 \frac{(F_{\rm av})_{ij}r_j}{\varphi r_i}
 =
 \begin{pmatrix}
 0&1\\
 \varphi^{-2}&\varphi^{-1}
 \end{pmatrix}.}
 \label{gfp:b:eq:parry-matriz-phi}
\end{equation}
La medida estacionaria de vértices es proporcional a
\((1,\varphi^2)\), luego
\begin{equation}
 \boxed{
 \frac{\mu_{\rm ar}}{\mu_{\rm vol}}=\varphi^{-2}.}
 \label{gfp:b:eq:parry-razon-phi}
\end{equation}
\end{theorem}

\begin{proof}
La suma de la fila \(i\) de \(P\) es
\[
 \frac{(F_{\rm av}r)_i}{\varphi r_i}=1.
\]
Como \(F_{\rm av}\) es simétrica, los vectores propios izquierdo y derecho
coinciden. Los pesos estacionarios son proporcionales a sus productos
coordenada a coordenada, es decir, a \((1,\varphi^2)\). La razón entre los
dos pesos es \(\varphi^{-2}\).
\end{proof}

\begin{remark}[Tres medidas que no deben confundirse]
La frecuencia cronológica de la órbita primitiva es
\[
 \nu_{\rm cyc}(\mathscr H_{\rm ar})=\frac13,
 \qquad
 \nu_{\rm cyc}(\mathscr H_{\rm vol})=\frac23.
\]
La medida uniforme \(\tau_{468}\) vive sobre el catálogo de emisiones. La
medida de Parry vive sobre historias admisibles de longitud arbitraria. La
primera cuenta instantes, la segunda cuenta emisiones y la tercera pondera
cilindros dinámicos. Ninguna es una aproximación de las otras.
\end{remark}

\section{Carácter multiplicativo y conservación cilíndrica}

Para una ruta admisible \(\gamma:i\to j\) de longitud \(n\), definimos
\begin{equation}
 \chi_P(\gamma)=\varphi^n\frac{r_i}{r_j}.
 \label{gfp:b:eq:caracter-parry-phi}
\end{equation}

\begin{proposition}[Multiplicatividad por concatenación]
\label{gfp:b:prop:multiplicatividad-parry-phi}
Si \(\gamma_1:i\to j\) y \(\gamma_2:j\to k\), entonces
\[
 \chi_P(\gamma_2\circ\gamma_1)
 =\chi_P(\gamma_2)\chi_P(\gamma_1).
\]
Sobre una ruta cerrada de longitud \(n\),
\[
 \chi_P(\gamma)=\varphi^n.
\]
\end{proposition}

\begin{proof}
La longitud se suma y la coordenada intermedia telescopa:
\[
 \varphi^{n_1+n_2}\frac{r_i}{r_k}
 =\left(\varphi^{n_1}\frac{r_i}{r_j}\right)
  \left(\varphi^{n_2}\frac{r_j}{r_k}\right).
\]
Si \(i=k\), el cociente terminal es uno.
\end{proof}

Normalizamos el vector izquierdo como
\[
 \ell=\frac{r}{1+\varphi^2},
 \qquad
 \ell^{\mathsf T}r=1.
\]
La medida de un cilindro es
\[
 \mu[x_0\cdots x_n]
 =\ell_{x_0}r_{x_n}\varphi^{-n}.
\]
Por tanto,
\begin{equation}
 V_n(x)
 :=\frac{\mu[x_0]}{\mu[x_0\cdots x_n]}
 =\varphi^n\frac{r_{x_0}}{r_{x_n}}
 =\chi_P(x_0\cdots x_n).
 \label{gfp:b:eq:volumen-cilindrico-phi}
\end{equation}
Para una dimensión tipada \(D>0\), la escala
\[
 a_n^{(D)}(x)=V_n(x)^{1/D}
\]
satisface la ley exacta
\begin{equation}
 \boxed{
 \bigl(a_n^{(D)}(x)\bigr)^D
 \mu[x_0\cdots x_n]=\mu[x_0].}
 \label{gfp:b:eq:conservacion-cilindrica-phi}
\end{equation}

En dimensión tres y sobre retornos cerrados,
\[
 \frac{a_9}{a_0}=\varphi^3,
 \qquad
 \frac{a_{27}}{a_0}=\varphi^9,
 \qquad
 \frac{a_{108}}{a_0}=\varphi^{36}.
\]
Si \(A_k=a_{9k}\), entonces
\begin{equation}
 A_{k+1}-2A_k+A_{k-1}
 =A_{k-1}(\varphi^3-1)^2>0.
 \label{gfp:b:eq:convexidad-retornos-phi}
\end{equation}
Ésta es convexidad en profundidad de retorno; no se interpreta como una
magnitud cinemática sin una carta adicional de duración.

\section{Memoria cuadrática y rama contraída}

Los autovalores de \(F_{\rm av}\) son
\[
 \varphi,
 \qquad
 -\varphi^{-1}.
\]
La rama contraída cambia de orientación en cada iteración. Para conservar
ese signo antes de elevar al cuadrado, pasamos a
\(\operatorname{Sym}^2(\mathbb R^2)\). En la base
\((x^2,2xy,y^2)\),
\begin{equation}
 \operatorname{Sym}^2(F_{\rm av})=
 \begin{pmatrix}
 0&0&1\\
 0&1&2\\
 1&1&1
 \end{pmatrix}.
 \label{gfp:b:eq:sym2-Fav-phi}
\end{equation}

\begin{theorem}[Espectro de la memoria de segundo orden]
\label{gfp:b:thm:espectro-sym2-phi}
\[
 \det\bigl(tI-\operatorname{Sym}^2(F_{\rm av})\bigr)
 =(t+1)(t^2-3t+1),
\]
y
\begin{equation}
 \boxed{
 \operatorname{Spec}\bigl(\operatorname{Sym}^2(F_{\rm av})\bigr)
 =\{\varphi^2,-1,\varphi^{-2}\}.}
 \label{gfp:b:eq:espectro-sym2-phi}
\end{equation}
El único modo estable positivo es \(\varphi^{-2}\).
\end{theorem}

\begin{proof}
Los autovalores de una potencia simétrica de grado dos son los productos
\(\lambda_1^2\), \(\lambda_1\lambda_2\) y \(\lambda_2^2\). Con
\(\lambda_1=\varphi\) y \(\lambda_2=-\varphi^{-1}\), se obtienen
\(\varphi^2\), \(-1\) y \(\varphi^{-2}\). La primera excede uno, la
segunda tiene módulo uno y la tercera pertenece a \((0,1)\).
\end{proof}

Sobre una ruta cerrada con \(V_n=\varphi^n\), el modo estable satisface
\begin{equation}
 \boxed{M_n=\varphi^{-2n}=V_n^{-2}.}
 \label{gfp:b:eq:modo-estable-phi}
\end{equation}

\section{Realización hiperbólica normalizada de la autoescala: \(\exp(2\log\varphi\,J_h)=F_{\mathrm{av}}^2\)}

Sea
\[
 A=F_{\rm av}^{\,2}=
 \begin{pmatrix}1&1\\1&2\end{pmatrix},
 \qquad
 J_h=\frac{2A-3I}{\sqrt5}.
\]
Entonces
\[
 J_h^2=I,
 \qquad
 A=\frac32I+\frac{\sqrt5}{2}J_h.
\]
Como
\[
 \cosh(2\log\varphi)=\frac32,
 \qquad
 \sinh(2\log\varphi)=\frac{\sqrt5}{2},
\]
se obtiene
\begin{equation}
 \boxed{
 \exp(2\log\varphi\,J_h)=F_{\rm av}^{\,2}.}
 \label{gfp:b:eq:realizacion-hiperbolica-phi}
\end{equation}
La constante de autoescala es, por tanto, simultáneamente el autovalor de
Perron de la incidencia y el parámetro exponencial de su realización
hiperbólica normalizada.

\section{Caracterización, horquillas y prolongación}

La construcción finita determina el rayo positivo de \(F_{\rm av}\). La
caracterización analítica posterior es
\[
 x^2-x-1=0,
 \qquad x>0.
\]
Las horquillas \eqref{gfp:b:eq:horquilla-fibonacci-phi} tienen extremos racionales
y anchura explícita \eqref{gfp:b:eq:anchura-fibonacci-phi}. Para cada nivel
nonádico \(K\), elegimos el menor \(m=m_{\varphi}(K)\) cuya horquilla queda
contenida en un único subcilindro entre los \(729\) elementos de la fibra
ambiente. La minimalidad convierte la localización en una función, no en una
elección tácita.

Los cilindros seleccionados son compatibles con los truncamientos y sus
diámetros tienden a cero. Su intersección contiene un único real; el
\cref{gfp:b:thm:rayo-perron-phi} lo identifica con \(\varphi\). Los cocientes de
Fibonacci son, pues, certificados racionales de la localización, no datos
decimales empleados para decidir la semilla.

\Needspace{22\baselineskip}
\section{Certificado interno de unicidad, compatibilidad y convergencia de la autoescala de Perron}

\begin{remark}[Certificado y control de validez: Refutación de la biografía de autoescala]
La construcción queda refutada si ocurre cualquiera de los hechos
siguientes:
\begin{enumerate}
 \setlength{\itemsep}{0.15\baselineskip}
 \item el predicado finito selecciona más de una órbita o más de una emisión;
 \item la recurrencia \(L_0,L_0,L_1,L_1\) no produce
       \eqref{gfp:b:eq:w30-phi};
 \item el calendario no produce exactamente los tres pares usados en
       \eqref{gfp:b:eq:N3-phi};
 \item se confunde un múltiplo de conteo \(N_{9},N_{27},N_{108}\) con una
       potencia de \(N_3\);
 \item \(F_{\rm av}\) posee otro rayo propio positivo;
 \item falla \eqref{gfp:b:eq:potencias-Fav-fibonacci} o la identidad de Cassini;
 \item las horquillas racionales no son encajadas o no contraen su anchura;
 \item se identifica la frecuencia cronológica, la medida uniforme del
       catálogo o la medida de Parry;
 \item una tabla decimal interviene antes de construir el rayo de Perron;
 \item un cilindro seleccionado no trunca al cilindro anterior.
\end{enumerate}
\end{remark}
